% ============================================================
% CHAPTER 3: METHODS
% University of Turku, Faculty of Technology
% MSc Thesis - Microbiome-Based BMI Prediction
% ============================================================

\chapter{Methods}
\label{ch:methods}

This chapter describes the dataset, computational pipeline, and experimental design employed in this thesis. All code is publicly available in the accompanying GitHub repository.

\section{Dataset: The Metalog Consortium}

\subsection{Data Sources and Merging}

The analysis draws on two complementary data products from the Metalog project, a large-scale aggregation of human gut metagenomes:

\begin{itemize}
    \item \textbf{Metadata Table} (\texttt{human\_extended\_wide.tsv}): Contains sample identifiers and phenotypic variables including age, sex, weight, height, and calculated BMI.
    \item \textbf{Taxonomic Profiles} (\texttt{human\_metaphlan4\_species.tsv}): Species-level relative abundances generated by MetaPhlAn4 from shotgun metagenomic reads.
\end{itemize}

Tables were merged on the common \texttt{sample\_id} key using an inner join, retaining only samples present in both sources. The merged dataset comprises \textbf{17,903 human samples} spanning diverse geographic origins, dietary patterns, and health statuses.

\subsection{Target Variable}

The prediction target is Body Mass Index (BMI), defined as:
\begin{equation}
\text{BMI} = \frac{\text{Weight (kg)}}{\text{Height (m)}^2}
\end{equation}
BMI is treated as a continuous outcome for regression modeling. The distribution in the dataset ranges from 14.5 to 58.2 kg/m\textsuperscript{2} with a mean of 25.3 and standard deviation of 5.1.

\subsection{Data Leakage Prevention}

A critical design decision concerns the exclusion of certain predictor variables to prevent data leakage:

\begin{itemize}
    \item \textbf{Weight:} Excluded because BMI is derived directly from weight. Including weight would provide the model with the ``answer'' and yield artificially inflated performance.
    \item \textbf{Age and Sex:} Excluded to isolate the microbial signal. While age and sex are known correlates of both BMI and microbiome composition, including them would conflate demographic prediction with microbiome-based inference. This study aims to establish a \emph{microbiome-only baseline}---demonstrating that gut composition alone contains predictive information about host phenotype.
\end{itemize}

\section{Computational Pipeline}

\subsection{Workflow Management: Nextflow}

The analysis pipeline is implemented in Nextflow, a domain-specific language for data-driven workflows \cite{ditommaso2017nextflow}. Nextflow offers:

\begin{itemize}
    \item \textbf{Reproducibility:} Containerized execution ensures identical results across environments.
    \item \textbf{Portability:} The same code runs on local machines, cloud platforms, and HPC clusters without modification.
    \item \textbf{Parallelism:} Independent tasks (e.g., training different models) execute concurrently, reducing wall-clock time.
\end{itemize}

The pipeline comprises two main processes:

\begin{enumerate}
    \item \texttt{Preprocess}: Loads the merged dataset, applies filters, and prepares data for modeling.
    \item \texttt{TrainModel}: Trains a specified machine learning algorithm and outputs performance metrics.
\end{enumerate}

\subsection{Preprocessing}

Preprocessing is performed using the \texttt{mikropml} R package and proceeds as follows:

\subsubsection{1\% Prevalence Filter}

Bacterial species present in fewer than 1\% of samples (approximately 179 individuals) are removed. This filter eliminates rare taxa unlikely to generalize across populations and substantially reduces dimensionality:

\begin{center}
\textbf{Features: 10,701 $\rightarrow$ 1,799} (83\% reduction)
\end{center}

This aggressive filtering accelerates computation without sacrificing predictive signal, as rare taxa contribute primarily noise to ensemble models.

\subsubsection{Variance Filtering (Top 1000)}

To address the high dimensionality and computational resource constraints, a variance-based feature selection step was implemented. We computed the variance for all 1,799 prevalence-filtered features and retained only the \textbf{Top 1,000 features} with the highest variance. This step reduces noise by eliminating features that show little variation across the population, focusing the model on the most potentially discriminative taxa.

\subsubsection{Centering and Scaling}

The selected 1,000 features are centered (mean = 0) and scaled (standard deviation = 1). Scaling parameters are computed on the training set and applied to the test set to prevent leakage.

\section{Machine Learning Models}

Four algorithms were benchmarked, representing a spectrum from linear to highly non-linear approaches:

\subsection{Elastic Net (GLMNet)}

Elastic Net combines L1 (Lasso) and L2 (Ridge) regularization. It performs implicit feature selection and serves as a linear baseline \cite{zou2005regularization}.

\subsection{Random Forest}

Random Forest constructs an ensemble of $B$ decision trees. We trained the model using 5-fold cross-validation. This bagging strategy decorrelates trees, reducing variance and capturing non-linear interactions \cite{breiman2001random}.

\subsection{XGBoost}

Extreme Gradient Boosting builds trees sequentially, with each tree correcting residuals from the previous iteration. We utilized the \texttt{xgbTree} method in \texttt{caret}, which is highly efficient and robust \cite{chen2016xgboost}.

\subsection{Support Vector Machine (Linear)}

While SVM with a radial basis function (RBF) kernel is powerful, it scales poorly ($O(n^3)$) with sample size. To enable training on the full dataset (18,000 samples) within local hardware memory limits, we employed \textbf{Linear SVM} (\texttt{svmLinear}). This approach finds the optimal separating hyperplane in the feature space with significantly lower computational overhead ($O(n)$) compared to kernel methods.

\section{Evaluation Framework}

\subsection{Cross-Validation}

Model performance is estimated using $k$-fold cross-validation (default $k=5$). The dataset is partitioned into $k$ folds; for each fold, the model trains on $k-1$ folds and evaluates on the held-out fold. Reported metrics represent the mean across folds.

\subsection{Metrics}

For regression tasks, primary metrics include:

\begin{itemize}
    \item \textbf{R-squared ($R^2$):} The proportion of variance explained by the model:
    \begin{equation}
    R^2 = 1 - \frac{\sum_i (y_i - \hat{y}_i)^2}{\sum_i (y_i - \bar{y})^2}
    \end{equation}
    \item \textbf{Root Mean Squared Error (RMSE):} The standard deviation of residuals:
    \begin{equation}
    \text{RMSE} = \sqrt{\frac{1}{n} \sum_i (y_i - \hat{y}_i)^2}
    \end{equation}
\end{itemize}

\subsection{Feature Importance}

For Random Forest, feature importance is quantified via \emph{increase in node purity} (Gini importance). Features splitting nodes effectively---those yielding large reductions in squared error---receive higher importance scores. The top 20 features are reported to identify candidate bacterial taxa for downstream biological investigation.

\section{Saturation Analysis}

To assess the relationship between sample size and performance, a saturation (learning curve) analysis was conducted:

\begin{enumerate}
    \item Subsamples of size $n \in \{1000, 5000, 10000, 15000, 17903\}$ were drawn.
    \item For each subsample, Random Forest was trained with 2-fold cross-validation.
    \item $R^2$ and training time were recorded.
\end{enumerate}

This analysis addresses whether additional samples yield diminishing returns and informs decisions about computational resource allocation.

\section{Hardware and Software}

Experiments were conducted on a consumer-grade Apple MacBook with 8GB RAM. Computationally intensive models (Random Forest on full data) required up to 16 hours of wall-clock time. The Nextflow pipeline is designed for seamless transition to HPC environments for future scaling.

\textbf{Software Stack:}
\begin{itemize}
    \item Nextflow 23.10.0
    \item R 4.3.1 with packages: \texttt{mikropml}, \texttt{caret}, \texttt{randomForest}, \texttt{xgboost}
    \item Conda for environment management
\end{itemize}
